%!TEX root = ../chapter06.tex 
%----------------------------------------------------------------------------------------
%	
%---------------------------------------------------------------------------------------- 
%----------------------------------------------------------------------------------------
%	 
%----------------------------------------------------------------------------------------



\newpage 
\loadgeometry{widelayout}  
\pagestyle{scrheadings} % Print headers again  
\KOMAoptions{
	headwidth=textwithmarginpar,
	headsepline=true,
	headsepline= 0.5pt,
	footwidth=textwithmarginpar,
} 
\onehalfspacing
\subsection{Exercices}



\begin{exercise}\label{chapter06:exercice:vecteur:coordonneesgraphique}\hfill \\
	\begin{minipage}{1\columnwidth}\centering
		\begin{tikzpicture}[scale=0.75 			]
			\tkzSetUpPoint[size=3, fill=black]
			\tkzfctset{tan style/.style={->,>=latex, color=red, line width=1.2pt}}   
			\tkzSetUpLine[line width= 1.2pt, add=.5 and .5] 
			
			
			\tkzDefPoint(7,-4){A} 
			\tkzDefShiftPoint[A](2,2){A'} 
			\tkzDrawSegment[->,>=latex, line width=1.2](A,A') 
			\tkzLabelSegment[above left](A,A') {\scriptsize $\vv{AA'}$}
			\tkzLabelPoint[below](A){\scriptsize$A$}\tkzLabelPoint[right](A'){\scriptsize$A'$}
			
			\tkzDefPoint(1,-2){B} 
			\tkzDefShiftPoint[B](3,-1){B'} 
			\tkzDrawSegment[->,>=latex, line width=1.2](B,B') 
			\tkzLabelSegment[above](B,B') {\scriptsize $\vv{BB'}$}
			\tkzLabelPoint[below](B){\scriptsize$B$}\tkzLabelPoint[right](B'){\scriptsize$B'$}
			
			\tkzDefPoint(-7,8){C} 
			\tkzDefShiftPoint[C](7,-3){C'} 
			\tkzDrawSegment[->,>=latex, line width=1.2](C,C') 
			\tkzLabelSegment[above](C,C') {\scriptsize $\vv{CC'}$}
			\tkzLabelPoint[below](C){\scriptsize$C$}\tkzLabelPoint[right](C'){\scriptsize$C'$}
			
			\tkzDefPoint(6,8){D} 
			\tkzDefShiftPoint[D](-2,4){D'} 
			\tkzDrawSegment[->,>=latex, line width=1.2](D,D') 
			\tkzLabelSegment[right](D,D') {\scriptsize $\vv{DD'}$}
			\tkzLabelPoint[below](D){\scriptsize$D$}\tkzLabelPoint[right](D'){\scriptsize$D'$}
			
			\tkzDefPoint(8,10){E} 
			\tkzDefShiftPoint[E](6,1){E'} 
			\tkzDrawSegment[->,>=latex, line width=1.2](E,E') 
			\tkzLabelSegment[above](E,E') {\scriptsize $\vv{EE'}$}
			\tkzLabelPoint[below](E){\scriptsize$E$}\tkzLabelPoint[right](E'){\scriptsize$E'$}
			
			
			
			\tkzInit[xmin=-8, ymin=-6,  xmax=15,ymax=12]  
			\begin{scope}[dashed] 
				\tkzGrid 
			\end{scope}  
			\tkzDrawX[ above=5pt,  label=$x$, font=\scriptsize, line width=1.2pt ] %, style=dashed
			\tkzDrawY[ right =5pt, label=$y$, font=\scriptsize, line width=1.2pt]  %, style=dashed
			\tkzRep[color  = Maroon,line width = 2pt, posxlabel=below, posylabel=left, xnorm=1,ynorm=1, xlabel={}, ylabel={}]
			\tkzSetUpPoint[size=3, fill=black]
			\tkzfctset{tan style/.style={->,>=latex, color=red, line width=1.2pt}}  
			\tkzDefPoint(2,5){M}
			\tkzDefPoint(3,1){N}
			\tkzDefPoint(7,6){P} 
			\tkzDefShiftPoint[M](2,2){M1}
			\tkzDefShiftPoint[M](3,-1){M2}
			\tkzDefShiftPoint[M](7,-3){N1}  
			
			\tkzDrawPoints[color=black, size=2](M,N,P)
			\tkzLabelPoint[above](M){\scriptsize$M$} 
			\tkzLabelPoint[below](N){\scriptsize$N$}  
			\tkzLabelPoint[above right](P){\scriptsize$P$} 
			\filldraw 
			(0,0) circle (2pt) node[align=left, below left] {\scriptsize $O$} ;
			\filldraw 
			(1,0) circle (2pt) node[align=left, below] {\scriptsize $I$} ;
			\filldraw 
			(0,1) circle (2pt) node[align=left, left] {\scriptsize $J$} ; 
		\end{tikzpicture} 
	\end{minipage}
	
	Dans le repère orthonormé $(O; I,J)$ :  
	\begin{enumerate}[leftmargin=*, label=\arabic*)]
		\item Donner les coordonnées des vecteurs $\vv{AA'}\begin{pmatrix}
			\ldots \\ \ldots
		\end{pmatrix}$, $\vv{BB'}\begin{pmatrix}
			\ldots \\ \ldots
		\end{pmatrix}$, $\vv{CC'}\begin{pmatrix}
			\ldots \\ \ldots
		\end{pmatrix}$, $\vv{DD'}\begin{pmatrix}
			\ldots \\ \ldots
		\end{pmatrix}$, $\vv{EE'}\begin{pmatrix}
			\ldots \\ \ldots
		\end{pmatrix}$. 
		\item \begin{enumerate}[leftmargin=*, label=\alph*)]
			\item Placer le point $M'$ image de $M$ par la translation $\vv{AA'}$.
			\item Placer l'image $N'$ de $N$ par la translation de vecteur $\vv{BB'}$.
			\item Placer l'antécédent $P'$ de $P$ par la translation de vecteur $\vv{CC'}$. 
			\item Placer l'antécédent $Q'$ de $Q$ par la translation de vecteur $\vv{DD'}$.  
			\item Placer l'image $R$ de $O$ par la translation de vecteur $\vv{EE'}$.      
		\end{enumerate}  
		\item \begin{multicols}{2} 
			\begin{enumerate}[leftmargin=*, label=\alph*)]
				\item   Placer le point $S$ tel que $\vv{ES}\begin{pmatrix}
					0\\2
				\end{pmatrix}$.
				\item   Placer le point $T$ tel que $\vv{AT}\begin{pmatrix}
					4\\0
				\end{pmatrix}$.
				\item   Placer le point $U$ tel que $\vv{BU}\begin{pmatrix}
					4\\3
				\end{pmatrix}$.
				\item   Placer le point $V$ tel que $\vv{VO}\begin{pmatrix}
					4\\3
				\end{pmatrix}$.
				\item   Placer le point $W$ tel que $\vv{WD}\begin{pmatrix}
					4\\-2
				\end{pmatrix}$.
				\item   Placer le point $X$ tel que $\vv{XE}\begin{pmatrix}
					-3\\5
				\end{pmatrix}$.
			\end{enumerate}  
		\end{multicols}
	\end{enumerate} 
\end{exercise} 

\begin{exercise}\label{chapter06:exercice:vecteurscoordonnees}\hfill \\ 
	Calculer les coordonnées des vecteurs $\vv{AB}$, $\vv{BC}$ et $\vv{CA}$  dans les cas suivants :
	\begin{multicols}{2}
		\begin{enumerate}[leftmargin=*,label=\alph*)]
			\item $A(3;-8)$, $B(-8;-9)$ et $C(0;8)$.
			\item $A(-1;1)$, $B(4;-4)$ et $C(5;-2)$.
			\item $A(0;0)$, $B(-4;5)$ et $C(3;\sqrt{2})$.
			\item $A(0;3)$, $B(-5;0)$ et $C(2;-2)$.
		\end{enumerate}
	\end{multicols}
\end{exercise}
\begin{example}Trouver les coordonnées des points $M$ et $N$ du repère $(O;I,J)$ dans les cas suivants\hfill \smallskip
	
	\noindent\parbox{0.47\linewidth}{
		$M$ est l'image de $F(-2;3)$ par la translation de vecteur $\vv{AB}\begin{pmatrix}
			3\\4
		\end{pmatrix}$
	}
	\hfill 
	\parbox{0.47\linewidth}{
		$G(5;3)$ est l'image de $N$ par la translation de vecteur $\vv{CD}\begin{pmatrix}
			-3\\2
		\end{pmatrix}$
	}	
	
	\vfill 
	
\end{example}

\begin{exercise}\label{chapter06:exercice:translationcoordonnees}\hfill \\
	On considère un repère $(O;I,J)$. Soit les points $A(1;-3)$, $B(-3;-4)$ et $C(-7;-8)$. \\
	Trouver par le calcul les coordonnées du point $M(x;y)$ dans les cas suivants : 
	\begin{enumerate}[leftmargin=*,label=\alph*)] 
		\item $M$ est l'image de $O$ par la translation de $\vv{AB}\begin{pmatrix}
			\ldots \\ \ldots
		\end{pmatrix}$. 
		\item $M$ est l'image de $A$ par la translation de $\vv{BC}$.
		\item $M$ est l'image de $I$ par la translation de $\vv{CA}$.
		\item $J$ est l'image de $M(x;y)$ par la translation de $\vv{CA}$.
	\end{enumerate} 
\end{exercise} 

\begin{exercise}[Égalités de vecteurs, direction, sens et norme]\hfill
	
	\noindent 
	\parbox[position=t]{1.0\linewidth}{
		\centering
		\begin{tikzpicture}[scale=0.80]
			\tkzInit[xmin=-1 ,xmax=7,ymin=-1, ymax=6]\tkzGrid 
			\tkzDefPoint(0,0){A}
			\tkzDefPoint(2,1){B}
			\tkzDefPoint(4,2){C}
			\tkzDefPoint(6,3){D}
			\tkzDefPoint(5,5){E}
			\tkzDefPoint(1,3){F}
			\tkzDrawPolygon[line width=1.2 pt](A,B,C,D,E,F) 
			\tkzLabelPoint[below left](A){$A$}
			\tkzLabelPoint[below right](B){$B$}
			\tkzLabelPoint[below right](C){$C$}
			\tkzLabelPoint[below right](D){$D$}
			\tkzLabelPoint[above right](E){$E$}
			\tkzLabelPoint[above left](F){$F$}  
			\tkzDrawSegments[line width=1.2pt](F,B F,C C,E)
		\end{tikzpicture} 
	} 
	
	En utilisant les points de la figure, indiquer 
	\begin{enumerate}[label=\alph*), leftmargin=*]
		\item 3 vecteurs de même direction que $\vv{AF}$ : \dotfill
		\item 3 vecteurs d'origine $C$ et de même direction que $\vv{EF}$ : \dotfill
		\item 2 vecteurs d'origine $B$ de même direction et sens contraire à $\vv{EF}$ :\dotfill 
		\item 2 vecteurs de même norme que $\vv{AF}$ mais pas de même direction : \dotfill 
		\item un vecteur égal à $\vv{DE}=$ : \dotfill  
		\item tous les vecteurs égaux à $\vv{AB}$  : \dotfill   
		\item tous les vecteurs égaux à $\vv{FE}$  : \dotfill   
	\end{enumerate}
\end{exercise} 



%----------------------------------------------------------------------------------------
%	 
%----------------------------------------------------------------------------------------


\newpage 

\begin{proof}[solution de l'exercice \ref{chapter06:exercice:vecteur:coordonneesgraphique}]\hfill 
	
	\noindent	$\vv{AA'}\begin{pmatrix}
		2 \\ 2
	\end{pmatrix}$, $\vv{BB'}\begin{pmatrix}
		3 \\ -1
	\end{pmatrix}$, $\vv{CC'}\begin{pmatrix}
		7 \\ -3
	\end{pmatrix}$, $\vv{DD'}\begin{pmatrix}
		-2 \\ 4
	\end{pmatrix}$, $\vv{EE'}\begin{pmatrix}
		6 \\ 1
	\end{pmatrix}$
	
	\noindent	\begin{minipage}{1\columnwidth}\centering
		\begin{tikzpicture}[scale=0.70 			]
			\tkzSetUpPoint[size=3, fill=black]
			\tkzfctset{tan style/.style={->,>=latex, color=red, line width=1.2pt}}   
			\tkzSetUpLine[line width= 1.2pt, add=.5 and .5] 
			
			
			\tkzInit[xmin=-8, ymin=-6,  xmax=15,ymax=12]  
			\begin{scope}[dashed] 
				\tkzGrid 
			\end{scope}  
			\tkzDrawX[ above=5pt,  label=$x$, font=\scriptsize, line width=1.2pt ] %, style=dashed
			\tkzDrawY[ right =5pt, label=$y$, font=\scriptsize, line width=1.2pt]  %, style=dashed
			\tkzRep[color  = Maroon,line width = 2pt, posxlabel=below, posylabel=left, xnorm=1,ynorm=1, xlabel={}, ylabel={}] 
			
			
			
			\tkzDefPoint(7,-4){A} 
			\tkzDefShiftPoint[A](2,2){A'} 
			\tkzDrawSegment[->,>=latex, line width=1.2](A,A') 
			\tkzLabelSegment[above left](A,A') {\scriptsize $\vv{AA'}$}
			\tkzLabelPoint[below](A){\scriptsize$A$}\tkzLabelPoint[right](A'){\scriptsize$A'$}
			
			\tkzDefPoint(1,-2){B} 
			\tkzDefShiftPoint[B](3,-1){B'} 
			\tkzDrawSegment[->,>=latex, line width=1.2](B,B') 
			\tkzLabelSegment[above](B,B') {\scriptsize $\vv{BB'}$}
			\tkzLabelPoint[below](B){\scriptsize$B$}\tkzLabelPoint[right](B'){\scriptsize$B'$}
			
			\tkzDefPoint(-7,8){C} 
			\tkzDefShiftPoint[C](7,-3){C'} 
			\tkzDrawSegment[->,>=latex, line width=1.2](C,C') 
			\tkzLabelSegment[above](C,C') {\scriptsize $\vv{CC'}$}
			\tkzLabelPoint[below](C){\scriptsize$C$}\tkzLabelPoint[right](C'){\scriptsize$C'$}
			
			\tkzDefPoint(6,8){D} 
			\tkzDefShiftPoint[D](-2,4){D'} 
			\tkzDrawSegment[->,>=latex, line width=1.2](D,D') 
			\tkzLabelSegment[right](D,D') {\scriptsize $\vv{DD'}$}
			\tkzLabelPoint[below](D){\scriptsize$D$}\tkzLabelPoint[right](D'){\scriptsize$D'$}
			
			\tkzDefPoint(8,10){E} 
			\tkzDefShiftPoint[E](6,1){E'} 
			\tkzDrawSegment[->,>=latex, line width=1.2](E,E') 
			\tkzLabelSegment[above](E,E') {\scriptsize $\vv{EE'}$}
			\tkzLabelPoint[below](E){\scriptsize$E$}\tkzLabelPoint[right](E'){\scriptsize$E'$}
			
			
			
			
			\tkzDefPoint(2,5){M}
			\tkzDefPoint(3,1){N}
			\tkzDefPoint(7,6){P} 
			\tkzDefShiftPoint[M](2,2){M1}
			\tkzDrawSegment[->,>=latex, line width=1.2,blue](M,M1) 
			\tkzLabelPoint[right](M1){$M'$}
			
			\tkzDefShiftPoint[N](3,-1){N1}
			\tkzDrawSegment[->,>=latex, line width=1.2,blue](N,N1) 
			\tkzLabelPoint[right](N1){$N'$}
			
			\tkzDefShiftPoint[P](-7,+3){P1}
			\tkzDrawSegment[->,>=latex, line width=1.2,blue](P1,P) 
			\tkzLabelPoint[right](P1){$P'$}
			
			
			\tkzDefPoint(11,2){Q} 
			\tkzDrawPoint(Q)\tkzLabelPoint[above right](Q){$Q$}
			\tkzDefShiftPoint[Q](2,-4){Q1}
			\tkzDrawSegment[->,>=latex, line width=1.2,blue](Q1,Q) 
			\tkzLabelPoint[right](Q1){$Q'$}
			
			
			\tkzDefShiftPoint[E](0,2){S}
			\tkzDrawSegment[->,>=latex, line width=1.2,red](E,S) 
			\tkzDrawPoint(S)\tkzLabelPoint[above right](S){$S$}
			
			\tkzDefShiftPoint[A](4,0){T}
			\tkzDrawSegment[->,>=latex, line width=1.2,red](A,T) 
			\tkzDrawPoint(T)\tkzLabelPoint[above right](T){$T$}
			
			\tkzDefShiftPoint[B](4,3){U}
			\tkzDrawSegment[->,>=latex, line width=1.2,red](B,U) 
			\tkzDrawPoint(U)\tkzLabelPoint[above right](U){$U$}
			
			\tkzDefPoint(0,0){O}	
			\tkzDefShiftPoint[O](-4,-3){V}
			\tkzDrawSegment[->,>=latex, line width=1.2,red](O,V) 
			\tkzDrawPoint(V)\tkzLabelPoint[above left](V){$V$}
			
			
			\tkzDefShiftPoint[O](+6,1){R}
			\tkzDrawSegment[->,>=latex, line width=1.2,blue](O,R) 
			\tkzDrawPoint(R)\tkzLabelPoint[above left](R){$R$}
			
			
			\tkzDefShiftPoint[D](-4,+2){W}
			\tkzDrawSegment[->,>=latex, line width=1.2,red](W,D) 
			\tkzDrawPoint(W)\tkzLabelPoint[above left](W){$W$}
			
			
			\tkzDefShiftPoint[E](+3,-5){X}
			\tkzDrawSegment[->,>=latex, line width=1.2,red](X,E) 
			\tkzDrawPoint(X)\tkzLabelPoint[above left](X){$X$}
			
			
			
			%		\tkzDefShiftPoint[M](3,-1){M2}
			%		\tkzDefShiftPoint[M](7,-3){N1} 
			%		  	  	\tkzDrawSegment[->,>=latex, line width=1.2](M1,P) 
			%		  	  	\tkzDrawSegment[->,>=latex, line width=1.2](M,M2) 
			%		  	  	\tkzDrawSegment[->,>=latex, line width=1.2](M2,P) 
			%			  	\tkzDrawSegment[->,>=latex, line width=1.2](M,N1) 
			%			  	\tkzDrawSegment[->,>=latex, line width=1.2](N1,P) 
			%			  	\tkzDrawSegment[->,>=latex, line width=1.2](N,N1) 
			%		  	\tkzLabelSegment[above](M,M1){\small $\vv{a}$}
			%			  	\tkzLabelSegment[above](M1,P) {\small $\vv{b}$}
			%			  	\tkzLabelSegment[above](M,M2){\small $\vv{b}$}
			%			  	\tkzLabelSegment[above](M2,P) {\small $\vv{a}$}
			%		%	  	\tkzLabelSegment[below](M,N1)  {\small $\vv{c}$}
			%			  	\tkzLabelSegment[right](N1,P)  {\small $\vv{d}$}
			%			  	\tkzLabelSegment[above](N,N1)   {\small $\vv{e}$} 	  	  	
			
			
			
			\tkzDrawPoints[color=black, size=2](M,N,P)
			\tkzLabelPoint[above](M){\scriptsize$M$} 
			\tkzLabelPoint[below](N){\scriptsize$N$}  
			\tkzLabelPoint[above right](P){\scriptsize$P$} 
			\filldraw 
			(0,0) circle (2pt) node[align=left, below left] {\scriptsize $O$} ;
			\filldraw 
			(1,0) circle (2pt) node[align=left, below] {\scriptsize $I$} ;
			\filldraw 
			(0,1) circle (2pt) node[align=left, left] {\scriptsize $J$} ; 
		\end{tikzpicture} 
	\end{minipage} 
\end{proof}

\begin{multicols}{2}  
	\begin{proof}[solution de l'exercice \ref{chapter06:exercice:vecteurscoordonnees}]\hfill 
		\begin{enumerate}[leftmargin=*,label=\alph*)] 
			\item  $\vv{AB}\begin{pmatrix}
				-11\\-1
			\end{pmatrix}$;
			$\vv{BC}\begin{pmatrix}
				8\\17
			\end{pmatrix}$;
			$\vv{CA}\begin{pmatrix}
				3\\-16
			\end{pmatrix}$;
			\item  $\vv{AB}\begin{pmatrix}
				5\\-5
			\end{pmatrix}$;
			$\vv{BC}\begin{pmatrix}
				1\\2
			\end{pmatrix}$;
			$\vv{CA}\begin{pmatrix}
				-6\\3
			\end{pmatrix}$;
			\item  $\vv{AB}\begin{pmatrix}
				-4\\5
			\end{pmatrix}$;
			$\vv{BC}\begin{pmatrix}
				7\\\sqrt{2}-5
			\end{pmatrix}$;
			$\vv{CA}\begin{pmatrix}
				-3\\-\sqrt{2}
			\end{pmatrix}$;
			\item  $\vv{AB}\begin{pmatrix}
				-5\\-3
			\end{pmatrix}$;
			$\vv{BC}\begin{pmatrix}
				7\\-2
			\end{pmatrix}$;
			$\vv{CA}\begin{pmatrix}
				-2\\3+\sqrt{2}
			\end{pmatrix}$;
		\end{enumerate} 
	\end{proof}
	\begin{proof}[solution de l'exercice \ref{chapter06:exercice:translationcoordonnees}]\hfill 
		\begin{enumerate}[leftmargin=*,label=\alph*)] 
			\item  $\vv{AB}\begin{pmatrix}
				-4\\-1
			\end{pmatrix}$ et $M(-4;-1)$.
			\item  $\vv{BC}\begin{pmatrix}
				-4\\-4
			\end{pmatrix}$ et $M(-3;-7)$.
			\item  $\vv{CA}\begin{pmatrix}
				8\\5
			\end{pmatrix}$; $I(1;0)$ et $M(9;5)$.
			\item $\vv{CA}\begin{pmatrix}
				8\\5
			\end{pmatrix}$; $J(0;1)$ et $M(-8;-4)$.
		\end{enumerate} 
	\end{proof} 
\end{multicols}
